\chapter{Implementation}\label{ch:impl}
This chapter describes the implementation of the Veritas toolchain as a journey from source-level specifications to independently checked machine code. The implementation consists of a refinement-typed source language, a bidirectional frontend checker, a typed SSA intermediate representation, a DTAL backend, an independent verifier, a compact runtime, and a bare-metal unikernel target. The chapter first gives the repository and pipeline structure, then follows the path taken by safety evidence through the compiler and verifier. Finally, Section~\ref{sec:worked-example} traces a binary search program through the full pipeline.

\section{Repository overview}\label{sec:repo-overview}

The implementation is deliberately broad: it includes the source language, compiler pipeline, verifier, executable generator, runtime, and bare-metal target. The source code is organised as a single Rust crate. Table~\ref{tab:repo} gives the main implementation areas rather than every internal module.
\begin{table}[h]
  \centering
  \small
  \begin{tabularx}{0.96\textwidth}{>{\raggedright\arraybackslash}p{0.28\textwidth} r >{\raggedright\arraybackslash}X}
    \hline
    \textbf{Area}                         & \textbf{Lines}  & \textbf{Responsibility}                     \\
    \hline
    \texttt{src/common/}                       & 687             & Shared AST, typed AST, and type definitions \\
    \texttt{src/frontend/}                     & 7,856           & Lexer, parser, and bidirectional type checker \\
    \texttt{src/backend/}                      & 17,391          & TIR, DTAL, optimisation, allocation, encoding, ELF, runtime \\
    \texttt{src/verifier/}                     & 7,372           & Independent DTAL verifier                   \\
    \texttt{src/main.rs}, \texttt{src/pipeline.rs} & 1,013       & CLI and pipeline orchestration              \\
    Other crate and binary files               & 262             & Module glue and helper binary               \\
    \hline
    \textbf{Total (excl.\ tests)}              & \textbf{34,581} &                                          \\
    \hline
  \end{tabularx}
  \caption{Repository structure and measured non-test Rust line counts.}
  \label{tab:repo}
\end{table}

\noindent The compiler, verifier, runtime, direct encoder, and test suite were written from scratch for this project. Outside \texttt{src/}, the repository contains 85 positive example programs, 35 negative examples, 424 Rust unit tests, integration tests for example verification and DTAL tampering, and an evaluation harness. The main third-party crates are \texttt{z3}, \texttt{chumsky}, \texttt{im}, and \texttt{ariadne}; these support infrastructure rather than implementing the Veritas language, compiler, verifier, runtime, or encoder.

Veritas is exposed as a command-line compiler over \texttt{.veri} source files. The normal path type-checks the source program, lowers it through TIR and DTAL, optionally verifies the resulting DTAL, and emits a freestanding ELF executable. The most important user-facing modes are \texttt{--verify}, which performs independent DTAL verification before native output is accepted, \texttt{--verify-dtal}, which checks a standalone \texttt{.dtal} file without using the source compiler, and \texttt{--target-bare-metal}, which emits a Multiboot-compatible binary for QEMU. This interface matters for the trust story: the verifier is not just an internal pass hidden inside the compiler, but can be run directly on the low-level artefact that the compiler produced.

\subsection{Pipeline overview}
Figure~\ref{fig:pipeline} summarises the compilation pipeline. Each named representation is shown in sequence, and arrows indicate the transformations between them. Type annotations are preserved through TAST, TIR, and DTAL, and are no longer needed after DTAL verification and final encoding to x86-64.
\begin{figure}[h]
  \centering
  \begin{tikzpicture}[
      scale=0.9,
      transform shape,
      stage/.style={draw, rounded corners=2pt, minimum width=1.65cm, minimum height=0.75cm, align=center, font=\scriptsize},
      evidence/.style={stage, fill=blue!9},
      plain/.style={stage, fill=gray!8},
      verifier/.style={stage, fill=green!12, minimum width=2.25cm},
      trusted/.style={stage, fill=green!7, minimum width=1.85cm},
      pass/.style={font=\scriptsize, text=black!70, align=center},
      arrow/.style={-Latex, thick}
    ]
    \node[plain] (src) at (0,0) {Source\\\texttt{.veri}};
    \node[plain] (ast) at (2.0,0) {AST};
    \node[evidence] (tast) at (4.0,0) {TAST\\types};
    \node[evidence] (tir) at (6.0,0) {TIR\\SSA + facts};
    \node[evidence] (vdtal) at (8.15,0) {Virtual\\DTAL};
    \node[evidence] (pdtal) at (10.45,0) {Physical\\DTAL};
    \node[verifier] (verify) at (10.45,-1.55) {Independent\\verifier};
    \node[trusted] (enc) at (8.15,-1.55) {Direct\\encoder};
    \node[plain] (elf) at (6.0,-1.55) {x86-64\\ELF};

    \draw[arrow] (src) -- (ast);
    \draw[arrow] (ast) -- (tast);
    \draw[arrow] (tast) -- (tir);
    \draw[arrow] (tir) -- (vdtal);
    \draw[arrow] (vdtal) -- (pdtal);
    \draw[arrow] (pdtal) -- (verify);
    \draw[arrow] (verify) -- (enc);
    \draw[arrow] (enc) -- (elf);

    \draw[dashed, rounded corners=3pt]
      (-0.55,0.55) rectangle (11.25,-0.55);
    \node[font=\scriptsize, fill=white, inner sep=1pt] at (5.35,0.62) {untrusted compiler};

    \draw[dashed, rounded corners=3pt]
      (7.0,-2.05) rectangle (11.65,-1.03);
    \node[font=\scriptsize, fill=white, inner sep=1pt] at (9.35,-2.12) {trusted checking and encoding core};

    \node[font=\scriptsize, align=center, text=blue!55!black] at (7.2,0.98)
      {proof-relevant safety evidence carried forward};
  \end{tikzpicture}
  \caption{The Veritas compilation pipeline. Blue stages carry type and constraint evidence; verification occurs after register allocation.}
  \label{fig:pipeline}
\end{figure}

\section{Source language and frontend checker}\label{sec:source-language}

This section presents the syntax, typing rules, and frontend implementation of the Veritas source language. These define the first contract in the pipeline: any program accepted by the type checker satisfies the properties encoded in its types.

\subsection{Syntax}

Veritas is a first-order imperative language with refinement-carrying types, fixed-size arrays, references, contracts, loop invariants, and bounded quantified specifications. The main syntactic categories are:
{\small
\[
\renewcommand{\arraystretch}{1.09}
\begin{array}{@{}r@{\ }c@{\ }l@{\hspace{0.8em}}p{0.22\textwidth}@{}}
  T, \tau & ::= &
  \parbox[t]{0.63\textwidth}{
    $\textsf{unit} \,|\, \textsf{int} \,|\, \textsf{bool} \,|\, [T; N] \,|\, \textsf{\&}T \,|\, \textsf{\&mut } T \,|\, \textsf{int}(N) \,|\, \{ x:\textsf{int} \mid P \}$
  }
  & \textbf{(Types)} \\[2pt]
  c & ::= &
  \underline{n} \,|\, \textsf{true} \,|\, \textsf{false}
  & \textbf{(Base values)} \\[2pt]
  \textsf{aop} & ::= &
  \texttt{+} \,|\, \texttt{-} \,|\, \texttt{*} \,|\, \texttt{/} \,|\, \texttt{\%}
  \qquad \textsf{uaop} ::= \texttt{-}
  & \textbf{(Arith. ops)} \\[2pt]
  \textsf{cop} & ::= &
  \texttt{==} \,|\, \texttt{!=} \,|\, \texttt{<} \,|\, \texttt{<=} \,|\, \texttt{>} \,|\, \texttt{>=}
  & \textbf{(Comparison ops)} \\[2pt]
  \textsf{lop} & ::= &
  \texttt{\&\&} \,|\, \texttt{||} \,|\, \texttt{==>}
  \qquad \textsf{ulop} ::= \texttt{!}
  & \textbf{(Logical ops)} \\[2pt]
  e & ::= &
  \parbox[t]{0.63\textwidth}{
    $c \,|\, x \,|\, e \textsf{ aop } e \,|\, e \textsf{ cop } e \,|\, e \textsf{ lop } e \,|\, \textsf{ulop}\; e \,|\, \textsf{uaop}\; e \,|\, f(\overrightarrow{e}) \,|\, {}$\\
    $x[e_i] \,|\, \kw{if } e_\textsf{cond}\; \{B\}\; \kw{else } \{B'\} \,|\, [e_\textsf{val}; e_\textsf{size}]$
  }
  & \textbf{(Expressions)} \\[4pt]
  P & ::= &
  \parbox[t]{0.63\textwidth}{
    $e \textsf{ cop } e \,|\, P \textsf{ lop } P \,|\, \textsf{ulop}\; P \,|\, {}$\\
    $\kw{forall } x \kw{in } e_\textsf{start}..e_\textsf{end}\;\{P\} \,|\, \kw{exists } x \kw{in } e_\textsf{start}..e_\textsf{end}\;\{P\}$
  }
  & \textbf{(Propositions)} \\[4pt]
  S & ::= &
  \parbox[t]{0.63\textwidth}{
    $\kw{let } x:T = e; \,|\, \kw{let mut } x:T = e; \,|\, x = e; \,|\, x[e_i] = e; \,|\, \kw{return } e; \,|\, e; \,|\, {}$\\
    $\kw{if } e_\textsf{cond}\; \{B\}\; \kw{else } \{B'\} \,|\, \kw{if } e_\textsf{cond}\; \{B\} \,|\, {}$\\ 
    $\kw{for } x \kw{in } e_\textsf{start}..e_\textsf{end}\; \optclause{\textsf{invariant } P}\; \{B\} \,|\, \kw{while } e_\textsf{cond}\; \optclause{\textsf{invariant } P}\; \{B\}$
  }
  & \textbf{(Statements)} \\[4pt]
  D & ::= &
  \parbox[t]{0.63\textwidth}{
    $\kw{fn } f(x_1:T_1, \ldots)\; \rightarrow T_\textsf{ret}\; \optclause{\textsf{requires } P}\;\optclause{\textsf{ensures } P}\; \{B\}$
  }
  & \textbf{(Function defs)} \\[2pt]
  \textsf{Prog} & ::= &
  D^*
  & \textbf{(Program)} \\[2pt]
  B & ::= &
  S^*\; \optclause{e}
  & \textbf{(Blocks)}
\end{array}
\]
}

Here $n$ ranges over integer literals, $x$ and $f$ over identifiers, and $N$ over statically known integer expressions used for array sizes and singleton types. $P$ is the shared specification language for refinements, contracts, loop invariants, and bounded quantifiers.

\subsection{Typing contexts}\label{sec:typing-contexts}

The typing context is a tuple $(\phi, \Gamma, \Delta, \Sigma_F)$:

\begin{center}
\small
\begin{tabularx}{0.95\textwidth}{>{\raggedright\arraybackslash}p{0.12\textwidth} >{\raggedright\arraybackslash}X}
\toprule
\textbf{Part} & \textbf{Meaning} \\
\midrule
$\phi$ & A set of refinement propositions known to be true: assumptions from branch conditions, preconditions, and refinement types. \\
$\Gamma$ & An immutable variable environment mapping names to types ($x : \tau$). \\
$\Delta$ & A mutable variable environment mapping names to a \textit{current type} and a \textit{master type} ($x : (\tau_c, \tau_m)$). We write $(x: \tau_c) \in \Delta$ to extract the current type and $(x: \tau_m) \in \Delta$ for the master type. \\
$\Sigma_F$ & A global function signature environment, read-only after an initial pass. \\
\bottomrule
\end{tabularx}
\end{center}

\subsubsection{Master types}

The master type $\tau_m$ is set at declaration and never changes; the current type $\tau_c$ may be refined through assignments and control flow but must remain a subtype of $\tau_m$. Thus a mutable \texttt{let mut x: T = e} records $T$ as the master type, allows the current type to narrow, and checks every assignment against the master type. At control-flow joins, current types from each branch are merged, but the result must still remain below the master type. This prevents reassignment from invalidating refinements.

The interaction between master types and mutable variables is described in Section~\ref{sec:typechecker}.

\subsection{Context joining}\label{sec:context-joining}

At control-flow merge points (i.e., after \textbf{if-else} and loops), the typing contexts from each branch must be joined. The $\textsf{join}$ operation computes a least upper bound for each mutable variable's type and intersects the proposition sets:

\begin{center}
\small
\begin{tabularx}{0.96\textwidth}{>{\raggedright\arraybackslash}p{0.24\textwidth} >{\raggedright\arraybackslash}X}
\toprule
\textbf{Case} & \textbf{Join rule} \\
\midrule
Same singletons & $\textsf{int}(N) \sqcup \textsf{int}(N) = \textsf{int}(N)$. \\
Different singletons & $\textsf{int}(N_1) \sqcup \textsf{int}(N_2) = \textsf{int}$ where $N_1 \ne N_2$. \\
Refined types & $\{v \mid P\} \sqcup \{v \mid Q\} = \{v \mid P \vee Q\}$. \\
Singleton and refined & Promote the singleton to $\{v \mid v = N\}$, then disjoin with the refined type's predicate. \\
Refined and unrefined & $\{v \mid P\} \sqcup \textsf{int} = \textsf{int}$. \\
Arrays & $[\tau_1; N] \sqcup [\tau_2; N] = [\tau_1 \sqcup \tau_2; N]$. \\
Propositions ($\phi$) & Intersect, keeping only propositions provable in both branches. \\
\bottomrule
\end{tabularx}
\end{center}

\noindent When $\textsf{join}$ is used in expression rules (e.g., \textsc{T-If-Expr}), the same type-joining rules apply to the result types of both branches.

\subsection{Judgement forms}\label{sec:judgement-forms}

Throughout the typing rules, $\tau$ is used as the metavariable for types.

The type system uses the following bidirectional typing judgements:

\begin{center}
\small
\begin{tabularx}{0.96\textwidth}{>{\raggedright\arraybackslash}p{0.17\textwidth} >{\raggedright\arraybackslash}X}
\toprule
\textbf{Judgement} & \textbf{Meaning} \\
\midrule
Synthesis & $\phi; \Delta; \Gamma \vdash e \Rightarrow (\phi'; \Delta'; \tau)$ synthesises type $\tau$ for expression $e$, potentially updating the context. \\
Checking & $\phi; \Delta; \Gamma \vdash e \Leftarrow \tau : (\phi'; \Delta')$ checks that $e$ has type $\tau$. \\
Subtyping & $\phi \vdash \tau_1 <: \tau_2$ states that $\tau_1$ is a subtype of $\tau_2$ under assumptions $\phi$. \\
Provability & $\phi \vdash P$ states that proposition $P$ is valid under assumptions $\phi$, decided by SMT. \\
\bottomrule
\end{tabularx}
\end{center}

The checking and synthesis judgements are connected by subsumption: if an expression synthesises type $\tau_1$ and the expected type is $\tau_2$, checking succeeds when $\tau_1 <: \tau_2$.

\subsection{Subtyping rules}

The subtyping relation $\phi \vdash \tau_1 <: \tau_2$ is defined by structural rules (reflexivity, singleton-to-int widening, refined-to-int widening, array covariance, master type transparency; listed in full in Appendix~\ref{app:source-rules}) and SMT-backed refinement implication. In each SMT-backed case, the propositions in $\phi$ are asserted as assumptions, the goal is negated, and the solver is checked for unsatisfiability:

$$\frac{\phi \vdash P[N/x]}{\phi \vdash \textsf{int}(N) <: \{x: \textsf{int} \mid P\}} \quad (\textsc{Sub-Singleton-Refine})$$

\subsection{Expression typing rules}\label{sec:expr-typing-rules}

The expression rules are standard bidirectional rules with three Veritas-specific behaviours:

\begin{center}
\small
\begin{tabularx}{0.96\textwidth}{>{\raggedright\arraybackslash}p{0.25\textwidth} >{\raggedright\arraybackslash}X}
\toprule
\textbf{Construct} & \textbf{Typing behaviour} \\
\midrule
Arithmetic & Constant-folds singleton operands and otherwise synthesises a refinement relating the result to the operands. \\
Array indexing & Synthesises the index type, extracts its logical proposition, and sends the bounds obligation $0 \le i < N$ to SMT. \\
Conditionals & Checks each branch under the condition and its negation, then joins the resulting contexts and result types. \\
Function calls & Checks argument subtyping and discharges the callee precondition after substituting actual arguments for formals. \\
\bottomrule
\end{tabularx}
\end{center}

The full expression rules, including the auxiliary $\textsf{val}$ and $\textsf{prop}$ functions used to extract index values and carried propositions, are given in Appendix~\ref{app:source-rules}.

\subsection{Statement typing rules}\label{sec:stmt-typing-rules}

Statement checking updates the typing context according to the statement's effect. Assignments synthesise the new value's type, check it against the variable's master type, and update the current type in $\Delta$. Loops add range facts for the induction variable and check any invariant twice: once on entry, and once after the body with the next iteration's value substituted. The full statement rules are given in Appendix~\ref{app:source-rules}.

\subsection{Function contracts}

Programmers may declare preconditions ($\mathbf{requires}\ P$) and postconditions ($\mathbf{ensures}\ P$) on functions. At each call site, the type checker verifies the precondition with actual arguments substituted for formals; at each return, it checks the return type and verifies the postcondition with the returned expression substituted for $\mathtt{result}$.

\subsection{Quantified specifications}

The constraint language supports bounded universal and existential quantifiers, plus implication, in specification positions only. This is enough to express array-wide properties such as sortedness and elementwise non-negativity without making quantifiers part of runtime computation. A useful implementation feature is \textit{refinement lifting}: an array typed as $[\{v\!:\!\textsf{int} \mid v \ge 0\}; N]$ automatically contributes the quantified fact $\forall i \in [0, N).\; \mathit{arr}[i] \ge 0$ when the array is in scope.

\subsection{Frontend implementation}\label{sec:frontend}

\subsubsection{Lexing and parsing}

The lexer and parser are built using \texttt{chumsky}. The frontend uses two expression parsers: a full parser for runtime expressions, and a restricted parser for refinement type annotations. This keeps type-position syntax unambiguous while reusing the same arithmetic and comparison language for specifications.

\subsubsection{Type checker}\label{sec:typechecker}

The type checker is the most substantial frontend component. It implements the bidirectional typing rules described in Section~\ref{sec:source-language}, using Z3 as an SMT oracle.

\detailhead{Architecture}The type checker is structured around the \texttt{TypingContext} $(\phi, \Gamma, \Delta, \Sigma_F)$ defined in Section~\ref{sec:typing-contexts}. Contexts are immutable values: each judgement takes a context as input and returns an updated context as output. The \texttt{im} crate makes branch cloning cheap through persistent data structures.

The type checker operates in two passes. The first collects all function signatures into $\Sigma_F$, enabling mutual recursion and allowing functions to be defined in any order. The second type-checks each function body against its signature.

\pagebreak[3]
\detailhead{Bidirectional checking}The implementation provides an explicit synthesis combinator, \texttt{synth\_expr}. Integer literals produce singleton types ($\mathtt{int}(n)$), variable lookups return the type from $\Gamma$ or $\Delta$, and arithmetic operations use \texttt{join\_op}. Checking first synthesises a type, then verifies by SMT-backed subtyping that it conforms to the expected type.

Statement checking (\texttt{check\_stmts}) updates the typing context according to each statement's effect: \texttt{let} bindings extend $\Gamma$ or $\Delta$, assignments update the current type in $\Delta$, and loops are checked by proving that their invariants hold on entry and are preserved by the body, as described in the statement typing rules for loops in Section~\ref{sec:stmt-typing-rules}.

\begin{verbatim}
synth_expr(ctx, expr) -> Result<(typed_expr, ty), TypeError>
check_stmt(ctx, stmt) -> Result<(typed_stmt, ctx'), TypeError>
check_stmts(ctx, stmts) -> Result<(typed_stmts, ctx'), TypeError>
check_program(program) -> Result<typed_program, TypeError>
\end{verbatim}

\detailhead{Singleton propagation and refinement synthesis}When both operands of an arithmetic operation are singletons, the checker folds at the type level: $\mathtt{int}(a) \oplus \mathtt{int}(b) = \mathtt{int}(a \oplus b)$. Otherwise it constructs a fresh refinement $\{v\!:\!\mathtt{int} \mid v = e_L \oplus e_R\}$, giving downstream stages enough information to verify computed indices such as \texttt{lo + (hi - lo) / 2}.

\detailhead{SMT oracle}\label{sec:smt-oracle}The SMT oracle wraps Z3 and provides a single entry point: $\texttt{is\_provable}(\phi, P)$. It asserts the assumptions, asserts $\lnot P$, and checks for unsatisfiability. Integer and boolean expressions map to the corresponding Z3 sorts; arrays use Z3 arrays; bounded quantifiers are translated to implications over native Z3 quantifiers.

The oracle also performs the refinement lifting described above, so element-level array refinements remain available to quantified proof obligations.

\detailhead{Error reporting}Type errors are structured \texttt{TypeError} values carrying source spans, expected and actual types, and contextual information. A reporting module formats these with \texttt{ariadne}, including the assumptions and failed proof goal where relevant.

\section{Typed intermediate representation and preservation}\label{sec:middle-end}

\subsection{Typed Intermediate Representation (TIR)}

The TIR is a typed SSA (Static Single Assignment) intermediate representation. Each function is a control-flow graph of basic blocks; each block contains a sequence of three-address instructions followed by a terminator (jump, conditional branch, or return). Every instruction carries type annotations from the source-level type checker, and every virtual register is defined exactly once.

\begin{center}
\small
\begin{tabularx}{0.96\linewidth}{@{}p{0.24\linewidth}X@{}}
\toprule
\textbf{Design choice} & \textbf{Implementation role} \\
\midrule
SSA form & Liveness analysis for register allocation is straightforward, and each virtual register has a single definition, so its type annotation has an unambiguous meaning. \\
Phi nodes & Control-flow joins become explicit program points where refinements from different predecessors can be joined rather than discarded silently. \\
Existential phi constraints & Loop counters and refined mutable variables can retain symbolic facts such as $\exists v.\; \mathtt{int}(v) \;\mathbf{where}\; (v \ge \mathit{start} \wedge v < \mathit{end})$, without fixing a concrete value. \\
\bottomrule
\end{tabularx}
\end{center}

The existential constraint is the key extension beyond ordinary SSA. It avoids the information loss that would occur if loop-carried values were simply widened to their base type.

\subsubsection{Representative instruction forms}

TIR is designed to sit close to both the typed source language and the DTAL backend. The implementation's instruction set includes ownership and region operations, but the core forms relevant to the compilation pipeline are summarised in Figure~\ref{fig:tir-syntax}.

In concrete terms, a TIR function stores its parameters, pre/postconditions, and a map of basic blocks; each block contains entry phi nodes, ordinary instructions, a terminator, and an entry register state. This makes the representation explicit enough for code generation to be largely structural, while still retaining the source-level type information needed for later verification.

\subsection{Lowering from TAST to TIR}

The lowering pass translates the typed AST to TIR. A \texttt{LoweringContext} maps source variables to virtual registers, tracks the current basic block, lowers statements in order, lowers any trailing expression to a result register, and packages the resulting blocks, phi nodes, entry states, and translated contracts into a \texttt{TirFunction}.

\subsubsection{Syntax-directed translation}

Lowering is syntax-directed: each typed source construct becomes a small SSA fragment while threading the variable-to-register map. Variable uses read the current register; \texttt{let} creates a fresh register; assignment updates the map so later uses see the new SSA value. Integer literals become \texttt{LoadImm}; arithmetic and comparisons become \texttt{BinOp} or \texttt{UnaryOp}; array operations become \texttt{ArrayLoad} and \texttt{ArrayStore} with explicit bounds constraints; calls carry argument registers, passing modes, ownership transfer, and result types. Conditionals introduce branch and merge blocks; loops introduce entry, header, body, and exit blocks with header phi nodes.

\begin{figure}[H]
\centering
\begin{verbatim}
block       ::= phi_node* instruction* terminator

phi_node    ::= PhiNode { dst, ty, incoming, existential_constraint? }

instruction ::= LoadImm { dst, value, ty }
              | Copy { dst, src, ty }
              | BinOp { dst, op, lhs, rhs, ty }
              | UnaryOp { dst, op, operand, ty }
              | ArrayLoad { dst, base, index, element_ty, 
                            bounds_constraint }
              | ArrayStore { base, index, value,
                             bounds_constraint }
              | Call { dst?, func, args, arg_types, arg_kinds, 
                       ownership, result_ty }
              | AssumeConstraint { constraint }
              | AssertConstraint { constraint }

terminator  ::= Jump { target }
              | Branch { cond, true_target, false_target,
                         true_constraint, false_constraint }
              | Return { value?, ownership }
\end{verbatim}
\caption{Representative TIR forms as defined in the implementation. The full implementation also includes ownership-transfer, borrowing, region, and allocation instructions.}
\label{fig:tir-syntax}
\end{figure}

\subsubsection{Ownership and borrowing through the pipeline}

Ownership is not treated as only a frontend check. The frontend tracks moved bindings, active borrows for each owner, reverse links from borrow variables back to their owners, and lexical borrow scopes. Creating \texttt{\&x} records a shared borrow; creating \texttt{\&mut x} requires no live borrow of the owner; moving or dropping an owner is rejected while conflicting borrows are still live.

\begin{center}
\small
\begin{tabularx}{\linewidth}{@{}p{0.19\linewidth}X@{}}
\toprule
\textbf{Stage} & \textbf{Ownership evidence kept explicit} \\
\midrule
Source checker & Binding state records owners, moved values, active shared borrows, active mutable borrows, and borrow scopes. \\
TIR & \texttt{MoveOwned}, \texttt{DropOwned}, \texttt{BorrowShared}, \texttt{BorrowMut}, and \texttt{BorrowEnd} make ownership changes visible before register allocation. \\
DTAL & \texttt{move\_owned}, \texttt{alias\_borrow}, \texttt{borrow\_mut}, \texttt{borrow\_end}, and \texttt{drop\_owned} expose the same events at the typed assembly level. \\
Verifier & Object identities, owned registers, and live borrows are checked after register allocation, so an incorrect backend cannot silently duplicate or alias ownership. \\
\bottomrule
\end{tabularx}
\end{center}

This gives ownership a low-level trust boundary. The verifier rejects duplicated owners, mutable borrows while aliases are live, and moves or drops while an active borrow remains. These checks run over allocated DTAL, at the same point as array-bounds and contract checks.

\subsubsection{Control flow}

\textbf{Conditionals} lower in the standard SSA style, except that each branch edge carries the logical fact derived from the condition. Thus \texttt{if x > 0} gives $x > 0$ on the taken edge and $\neg(x > 0)$ on the other. Phi nodes are inserted only where the assigned register differs.

\textbf{For-loops} use entry, header, body, and exit blocks. The induction variable is a header phi node with an existential range constraint $[\mathit{start}, \mathit{end})$, and mutable variables in scope receive loop-carried phi nodes. \textbf{While-loops} use the same shape, but rely on the translated branch condition and any programmer-supplied invariant rather than an explicit interval constraint.

\subsubsection{Constraint propagation}

Branch conditions are translated into the TIR constraint language: arithmetic comparisons, logical connectives, bounded quantifiers, and array indexing, but over virtual registers rather than source variables. A substitution pass rewrites names such as \texttt{"lo"} to registers such as \texttt{"v11"} before emitting constraints. Loop invariants become \texttt{AssertConstraint} instructions, later translated to DTAL \texttt{assert}s, so the verifier re-checks the proof-relevant facts independently.

\section{DTAL backend and physical code generation}\label{sec:backend}

\subsection{DTAL target language}\label{sec:dtal-target-lang}

As introduced in Section~\ref{sec:dtal-background}, Veritas targets a first-order variant of Xi and Harper's DTAL design~\cite{Xi2001DTAL}. In their original work, Xi and Harper formally established the soundness of DTAL, proving that well-typed programs cannot violate memory safety or the constraints of the type system. Veritas builds on this theoretical foundation: the compiler emits DTAL annotated with type information, and a standalone verifier re-derives the types from scratch. If the compiler introduces an error, the annotations will be inconsistent and the verifier will reject the program. This section describes the concrete instruction set and typing judgements of the Veritas DTAL.

\subsubsection{Instruction set}\label{sec:dtal-instr-set}

The DTAL instruction set resembles a standard RISC assembly language but is augmented with type annotations and constraint assertions. It operates over a set of virtual registers (which are later allocated to x86-64 physical registers) and linear memory:

\begin{verbatim}
data      ::= mov rd, rs | mov rd, n | load rd, [rb, ri]
            | store [rb, ri], rs
arith     ::= binop rd, rs1, rs2 | addimm rd, rs, n | not rd, rs
compare   ::= cmp r1, r2 | cmp r1, n | setcc rd, cond
control   ::= jmp L | branch cond L | call L | ret
stack     ::= push rs | pop rd | alloca rd, n
annot     ::= type r: tau | assert C
\end{verbatim}

\noindent The annotation instructions are the distinguishing feature of DTAL. Every instruction carries a type annotation on its destination register, and \texttt{assert} instructions embed constraint assertions (e.g., bounds proofs, loop invariants) directly into the instruction stream. These annotations serve as verification evidence: the compiler emits them, and the verifier checks them independently.

\subsubsection{DTAL type system}

The DTAL type system largely mirrors the source-level type system: it retains base types, singleton integers $\textsf{int}(N)$, refinement types, and arrays. The main difference is that these types classify \emph{registers} rather than source variables, so the indices appearing in types are low-level expressions over register contents. In particular, DTAL index expressions may mention constants, register names, arithmetic ($+, -, \times, \div, \%$), and $\textsf{select}(\mathit{arr}, i)$ for array element access.

The most important addition is the use of \textit{existential types} to represent results that are known only symbolically at the assembly level. When arithmetic on non-singleton operands cannot preserve an exact singleton type, the result is represented as an existential such as $\exists x.\; \textsf{int}(x) \;\mathbf{where}\; C$. This allows the verifier to retain enough logical structure for later proof obligations without requiring every intermediate value to have a concrete singleton type.

To make this precise, DTAL index expressions are given by:
\[
\begin{aligned}
  e_i ::=\, & n \,|\, r \,|\, e_i + e_i \,|\, e_i - e_i \,|\, e_i \times e_i \,|\, e_i \div e_i \,|\, e_i \bmod e_i \,|\, \textsf{select}(\mathit{arr}, e_i)
\end{aligned}
\]

\noindent where $n$ ranges over integer constants and $r$ over registers. The DTAL constraint language preserves the same logical fragment used in source-level specifications: comparisons, logical connectives, and bounded quantifiers. However, it is represented separately as a first-order language over DTAL index expressions rather than directly as source expressions. These constraints are threaded through the verifier's type state and refined by comparisons, branches, assertions, and function contracts.

\subsubsection{DTAL typing judgements}

DTAL type checking operates over a \textit{type state} $(\Gamma, \Phi)$ where $\Gamma$ maps registers to their current types and $\Phi$ is the constraint context at the current program point. The core judgement is $\Phi; \Gamma \vdash \textsf{instr} \Rightarrow \Gamma'$. The representative rules are straightforward: immediates produce singleton types; arithmetic over singleton operands preserves exact symbolic values; arithmetic over refined or existential operands produces an existential result; loads and stores require an SMT proof that the index is within bounds; and conditional branches refine $\Phi$ differently on taken and fall-through edges. Stores also version arrays so the SMT encoding can distinguish values before and after mutation. The complete rule set is given in Appendix~\ref{app:dtal-rules}.

\subsection{Backend implementation}

The code generator translates the TIR forms of Figure~\ref{fig:tir-syntax} to DTAL, the dependently typed assembly language described in Section~\ref{sec:dtal-target-lang}. At this stage DTAL still uses virtual registers, but each basic block is annotated with a \textit{type state}: a mapping from registers to types together with the constraints known to hold at that program point. These annotations are the core of the type-preserving design, because they give the standalone verifier enough information to re-check the program without access to the source.

Most TIR instructions lower directly to DTAL. Phi nodes are implemented as parallel copies on predecessor edges, and array operations become address computations plus loads or stores with the relevant bounds constraints attached.

\begin{center}
\small
\begin{tabularx}{0.96\linewidth}{@{}p{0.24\linewidth}X@{}}
\toprule
\textbf{Backend stage} & \textbf{Engineering work} \\
\midrule
DTAL generation & Emits block entry states, instruction types, assertions, and translated contracts, so the verifier can re-check the low-level program without seeing the source. \\
Optimisation & Runs constant folding, peephole rewrites, copy propagation, dead-code elimination, LICM, and load-op fusion to a bounded fixed point. Rewritten instructions continue to carry verifier-visible type evidence. \\
LICM & Uses dominators and natural-loop detection, and only hoists computations whose operands and memory effects make the move safe. Constraint assertions are treated as visible effects and are never eliminated. \\
Register allocation & Uses Chaitin--Briggs graph colouring over 11 allocatable x86-64 registers, with a linear-scan allocator implemented for comparison. Values live across calls are constrained to callee-saved registers. \\
Calling convention & Handles prologue and epilogue generation, caller-saved spills, argument-move cycles, and target-specific expansions such as division. \\
\bottomrule
\end{tabularx}
\end{center}

These stages are deliberately below the source-language safety proof. Optimisation, allocation, spill insertion, ABI lowering, and instruction expansion are untrusted transformations whose output is checked at the physical-DTAL boundary.

\subsubsection{Physical allocation}\label{sec:physalloc}

The register allocator produces a virtual-to-physical mapping; the \textit{physical allocation pass} applies it to produce DTAL over physical registers. This is the central untrusted transformation in the default pipeline: its output is verified before code generation proceeds. DTAL registers map directly to x86-64 registers:

\begin{table}[h]
\centering
\small
\begin{tabular}{l l l l}
\hline
\textbf{DTAL} & \textbf{x86-64} & \textbf{Role} & \textbf{Saved by} \\
\hline
\texttt{r0}--\texttt{r5} & \texttt{rdi}, \texttt{rsi}, \texttt{rdx}, \texttt{rcx}, \texttt{r8}, \texttt{r9} & Arguments 1--6 & Caller \\
\texttt{r6}--\texttt{r7} & \texttt{r10}, \texttt{r11} & Scratch & Caller \\
\texttt{r8}--\texttt{r12} & \texttt{rbx}, \texttt{r12}--\texttt{r15} & General & Callee \\
\texttt{lr} & \texttt{rax} & Return value & Caller \\
\texttt{sp} & \texttt{rsp} & Stack pointer & none \\
\texttt{fp} & \texttt{rbp} & Frame pointer & none \\
\hline
\end{tabular}
\caption{DTAL physical register mapping to x86-64.}
\label{tab:regmap}
\end{table}

Beyond register renaming, physical allocation emits prologues and epilogues, expands division into the x86-64 \texttt{rax}/\texttt{rdx} protocol, handles caller-saved spills, breaks argument-move cycles with a scratch register, and remaps all constraints from virtual to physical registers so that the verifier checks the actual allocation result.

\subsubsection{Direct encoding}

The \textit{direct encoder} translates physically allocated DTAL to x86-64 machine code without using an external assembler. Because physical allocation has already performed division expansion, calling-convention lowering, and prologue generation, the encoder is close to a 1:1 mapping and is only 328 lines. This split replaces an earlier monolithic trusted x86-64 lowering pass with an untrusted physical allocation pass plus a small trusted encoder, which is the intended TCB reduction.

\subsubsection{Binary encoding and ELF generation}

The direct encoder emits x86-64 machine code from physical DTAL, using a two-pass scheme to resolve label offsets. The resulting byte stream is then wrapped in a minimal ELF executable with entry point \texttt{main}; for the Linux target this is an ELF64 binary, while the bare-metal target uses a separate Multiboot-compatible layout described below.

\section{Independent verifier}\label{sec:verifier}

A key contribution of the type-preserving architecture is that the DTAL output can be verified independently of the compiler. The verifier (7,372 lines) re-checks the DTAL program from scratch, using only the type annotations embedded in the assembly. It follows the derivation-based verification approach of Xi and Harper~\cite{Xi2001DTAL}: types are \textit{re-derived} from instruction semantics, not trusted from compiler annotations.

\subsection{Verification algorithm}

The verifier processes each function independently. Its core loop is deliberately derivational: compiler annotations are cross-checks, not trusted facts.

\begin{center}
\small
\begin{tabularx}{0.96\linewidth}{@{}p{0.24\linewidth}X@{}}
\toprule
\textbf{Verifier step} & \textbf{Check performed} \\
\midrule
Initialise block state & Load the block's declared register types and constraints. \\
Execute instruction & Re-derive result types from instruction semantics, e.g. \texttt{mov rd, 5} gives $\texttt{rd} : \mathtt{int}(5)$ and arithmetic derives symbolic result types. \\
Check annotation & Compare the derived type with the compiler-supplied type annotation. \\
Discharge assertion & Prove each \texttt{.assert} using the verifier's own SMT oracle. \\
Coerce exit state & Check that jumps and branches satisfy the successor block's declared entry state. \\
\bottomrule
\end{tabularx}
\end{center}

\subsection{State coercion at control-flow edges}

State coercion checks compatibility across block boundaries. When block A jumps to block B, A's exit state must be a subtype of B's declared entry state: each required register type must be satisfied, and each declared successor constraint must follow from A's exit constraints. For conditional branches, the verifier first refines the two exits, so \texttt{cmp r1, r2} followed by \texttt{jl target} contributes $\texttt{r1} < \texttt{r2}$ on the taken edge and $\texttt{r1} \ge \texttt{r2}$ on the fall-through edge.

\subsection{Array bounds verification}

Array loads and stores are the most critical verifier instructions. For a \texttt{load rd, [base + offset]}, the verifier checks that \texttt{base} has array type $[\tau; N]$, proves $0 \le \texttt{offset} < N$, derives element type $\tau$ for \texttt{rd}, and emits a select fact linking the loaded value to the array's logical representation. Stores additionally emit:
\[
  \mathit{arr}_{\mathit{new}}[\texttt{offset}] = \texttt{value}
  \qquad
  \forall k \ne \texttt{offset}.\; \mathit{arr}_{\mathit{new}}[k] = \mathit{arr}_{\mathit{old}}[k].
\]
The array is versioned on each store, so SMT can distinguish array states before and after mutation.

\subsection{Interprocedural contract verification}

For calls, the verifier proves the callee precondition after substituting actual argument registers for formal parameters, assigns the declared return type to the destination register, and adds the substituted postcondition to the caller's constraint set. This enables modular verification: each function is checked in isolation, relying on contracts rather than inlining.

\subsection{SMT oracle}

The verifier uses its own Z3 instance, independent of the compiler's oracle. It follows the same proof-by-refutation structure, but first applies a fast syntactic check for common entailments before invoking SMT. This matters because the verifier issues many small queries, often one per instruction.

\subsection{Verification after register allocation}

The default physical pipeline verifies \textit{after} register allocation has mapped virtual registers to physical x86-64 registers. This makes register allocation, spill insertion, and calling-convention lowering untrusted: bugs should appear as type derivation mismatches or unprovable constraints in the verified output. A legacy mode can verify virtual-register DTAL earlier, but that leaves more backend work trusted.

This was one of the most important implementation shifts in the project. Verifying virtual DTAL is easier, because registers are still compiler-invented temporaries and the code is close to the typed intermediate representation. It is also less compelling as a trust boundary, because a later allocator or instruction-lowering pass could still corrupt the executable. Physical verification moves the check below those passes. The price is that the verifier must understand real backend artefacts: scratch registers, spill slots, calling convention obligations, prologues, epilogues, and instruction expansions such as division. The payoff is that a larger and more failure-prone part of the compiler is outside the TCB.

\subsection{Trust model}

The trust model and TCB reduction phases are described in Section~\ref{sec:design-decisions}. The verification architecture places the trust boundary between the physical DTAL verifier and the direct encoder: the entire compiler (frontend, middle end, backend, register allocator) is untrusted, and any bug in these components manifests as a verification failure rather than a safety violation.

\section{Runtime and bare-metal target}\label{sec:runtime}

The compiler links a small hand-assembled x86-64 runtime into generated executables. The 418-byte runtime provides \texttt{print\_int}, \texttt{print\_char}, and \texttt{read\_int} through Linux syscalls, and \texttt{port\_in}/\texttt{port\_out} through x86 I/O instructions. The type checker enforces which intrinsics are available on Linux and bare metal. The runtime also contains three internal region-management entry points used by hosted ownership, and remains part of the trusted computing base.

\subsection{Bare-metal target and unikernel}\label{sec:unikernel}

\texttt{--target-bare-metal} produces a Multiboot-compliant ELF32 that boots directly on x86-64 hardware or QEMU, removing the OS from the trusted computing base. The binary begins with an approximately 200-byte trusted bootstrap for Multiboot setup, paging, long-mode entry, stack initialisation, the call to \texttt{main}, and the final halt loop. The bootstrap identity-maps the first 8\,MB using 2\,MB pages, enables PAE and long mode, loads a minimal GDT, performs the far jump into 64-bit code, and then hands control to verified Veritas code. It remains trusted, but its size makes it a plausible audit target rather than another compiler-scale artefact.

To demonstrate end-to-end verified I/O on bare metal, the project includes a UART 16550 serial driver written entirely in Veritas. It initialises the serial port at COM1, configures the baud rate and FIFO, waits on the Line Status Register before each write, writes bytes through \texttt{port\_out}, and prints integers recursively. The driver is type-checked and compiled through the verified DTAL pipeline, producing a QEMU-bootable serial-console binary.

\begin{figure}[H]
  \centering
  \begin{tikzpicture}[
      scale=0.76,
      transform shape,
      layer/.style={draw, rounded corners=2pt, minimum width=11.0cm, minimum height=0.72cm, align=center, font=\scriptsize},
      trusted/.style={layer, fill=green!10},
      verified/.style={layer, fill=blue!9},
      hardware/.style={layer, fill=gray!10},
      arrow/.style={-Latex, thick}
    ]
    \node[hardware] (hw) at (0,0) {x86-64 hardware: CPU, RAM, UART 16550 serial port, halt loop};
    \node[trusted] (boot) at (0,0.95) {trusted bootstrap: Multiboot header, paging, long-mode entry, stack setup, call \texttt{main}};
    \node[trusted] (rt) at (0,1.9) {small trusted substrate: direct encoder output + port I/O intrinsics};
    \node[verified] (drv) at (0,2.85) {verified Veritas serial driver: \texttt{serial\_init}, \texttt{serial\_write}, \texttt{serial\_print\_int}};
    \node[verified] (app) at (0,3.8) {verified Veritas application code compiled through physical DTAL};

    \draw[arrow] (app) -- (drv);
    \draw[arrow] (drv) -- (rt);
    \draw[arrow] (rt) -- (boot);
    \draw[arrow] (boot) -- (hw);

    \draw[dashed, rounded corners=3pt] (-5.85,2.45) rectangle (5.85,4.25);
    \node[font=\scriptsize, fill=white, inner sep=1pt] at (0,4.36) {checked by the Veritas frontend and physical-DTAL verifier};

    \draw[dashed, rounded corners=3pt] (-5.85,0.55) rectangle (5.85,2.25);
    \node[font=\scriptsize, fill=white, inner sep=1pt] at (0,0.43) {remaining trusted code is small enough to audit directly};
  \end{tikzpicture}
  \caption{Bare-metal Veritas binary structure. The operating system is removed from the trusted computing base; the application and serial driver are compiled from Veritas.}
  \label{fig:bare-metal-unikernel}
\end{figure}

\section{Worked example: binary search}\label{sec:worked-example}
The preceding sections describe each pipeline stage in isolation. This section demonstrates the full pipeline by tracing a \texttt{binary\_search} program, from source to verified DTAL. Binary search features most of the verification and constraint features of Veritas: quantified preconditions, refinement types, loop invariants, array bounds proofs, existential types, branch constraint refinement, and symbolic arithmetic.

\begin{center}
\small
\begin{tabularx}{0.96\linewidth}{@{}p{0.22\linewidth}X@{}}
\toprule
\textbf{Pipeline point} & \textbf{What the example shows} \\
\midrule
Source & The sortedness precondition, loop invariant, and midpoint access state the safety argument. \\
Frontend & The checker turns those facts into SMT obligations, especially the proof that \texttt{arr[mid]} is in bounds. \\
TIR & SSA registers and existential phi constraints keep the loop facts explicit. \\
Physical DTAL & Register allocation and x86-specific expansion rewrite the program before the independent verifier accepts it. \\
\bottomrule
\end{tabularx}
\end{center}

\subsection{Source language}
The source program searches a sorted 10-element array for a target value, returning its index or $-1$ if not found. The proof-relevant fragment is:
\begin{verbatim}
fn binary_search(arr: [int; 10], target: int) -> int
    requires forall i in 0..9 { arr[i] <= arr[i + 1] }
{
    let mut lo: {v: int | v >= 0} = 0;
    let mut hi: int = 9;
    let mut result: int = -1;

    for iter in 0..20
        invariant lo >= 0 && hi < 10
    {
        if lo <= hi {
            let mid: int = lo + (hi - lo) / 2;
            let val: int = arr[mid];
            ... update lo, hi, and result ...
        } else {
            ...
        }
    }

    result
}
\end{verbatim}
Several features are worth highlighting:

\begin{center}
\small
\begin{tabularx}{0.96\textwidth}{>{\raggedright\arraybackslash}p{0.32\textwidth} >{\raggedright\arraybackslash}X}
\toprule
\textbf{Source feature} & \textbf{Role in the proof} \\
\midrule
\texttt{requires forall i in 0..9 \{ arr[i] <= arr[i+1] \}} & Quantified precondition stating sortedness, checked at every call site. \\
\texttt{\{v: int | v >= 0\}} & Master refinement on \texttt{lo}; every assignment must preserve non-negativity. \\
\texttt{invariant lo >= 0 \&\& hi < 10} & Loop invariant bounding both endpoints, checked initially and after the body. \\
\texttt{arr[mid]} & Critical access; the checker must prove $0 \le \mathtt{mid} < 10$. \\
\texttt{for iter in 0..20} & Bounded iteration counter, later preserved as an SSA and DTAL range fact. \\
\bottomrule
\end{tabularx}
\end{center}

\subsection{Frontend type checking}

The type checker discharges three key proof obligations for this function. All are reduced to SMT queries of the form: assert the current context $\phi$, assert $\lnot G$ (the negated goal), and check for unsatisfiability.

\begin{center}
\small
\begin{tabularx}{0.96\textwidth}{>{\raggedright\arraybackslash}p{0.26\textwidth} >{\raggedright\arraybackslash}X}
\toprule
\textbf{Obligation} & \textbf{Why it matters} \\
\midrule
Invariant base case & Shows the loop starts in a state satisfying $\mathtt{lo} \ge 0 \wedge \mathtt{hi} < 10$. \\
Invariant preservation & Shows every branch re-establishes the same invariant before the back edge. \\
Array bounds & Shows the computed midpoint is always a valid index before loading \texttt{arr[mid]}. \\
\bottomrule
\end{tabularx}
\end{center}

The base case follows from the singleton facts $\mathtt{lo}=0$ and $\mathtt{hi}=9$. Preservation is checked branch-by-branch: updates to \texttt{lo} preserve non-negativity, updates to \texttt{hi} preserve the upper bound, and context joining reconciles the branch results. The critical array access occurs under \texttt{lo <= hi}, giving the context:
\[
  \phi = \{\, \mathtt{lo} \ge 0, \;\; \mathtt{hi} < 10, \;\; \mathtt{lo} \le \mathtt{hi}, \;\; \mathtt{mid} = \mathtt{lo} + (\mathtt{hi} - \mathtt{lo}) / 2 \,\}
\]
The SMT oracle proves $0 \le \mathtt{mid} < 10$ from these facts, confirming the bounds proof for \texttt{arr[mid]}.

\subsection{Lowering to TIR (SSA)}
The lowering pass translates the function into a control-flow graph of 13 basic blocks. Figure~\ref{fig:bs-cfg} shows the structure, with blocks grouped by their role in the algorithm.
\begin{figure}[H]
  \centering
  \begin{tikzpicture}[
      scale=1.0, transform shape,
      block/.style={rectangle, draw, thick, text centered,
          minimum height=0.62cm, minimum width=0.98cm,
          font=\small\ttfamily},
      hotblock/.style={rectangle, draw=black, very thick, fill=yellow!18, text centered,
          minimum height=0.62cm, minimum width=0.98cm,
          font=\small\ttfamily},
      arrow/.style={-Stealth, thick},
      lbl/.style={font=\scriptsize, text=black!60},
      elbl/.style={font=\scriptsize, text=black!60, fill=white, inner sep=1pt}
    ]
    \node[block] (bb0) at (0,0) {bb0};
    \node[lbl, right=0.15cm of bb0] {entry};

    \node[block] (bb1) at (0,-1.6) {bb1};
    \node[lbl, left=0.15cm of bb1] {loop header};

    \node[block] (bb2) at (-2,-3.4) {bb2};
    \node[lbl, left=0.15cm of bb2] {body};
    \node[block] (bb3) at (3,-3.4) {bb3};
    \node[lbl, right=0.15cm of bb3] {exit};

    \node[hotblock] (bb4) at (-3.5,-5.2) {bb4};
    \node[block] (bb5) at (0.5,-5.2) {bb5};

    \node[block] (bb7) at (-5,-7.0) {bb7};
    \node[block] (bb8) at (-2,-7.0) {bb8};

    \node[block] (bb9) at (-5,-8.8) {bb9};
    \node[block] (bb10) at (-3,-8.8) {bb10};
    \node[block] (bb11) at (-1.7,-8.8) {bb11};

    \node[block] (bb12) at (-2.3,-10.4) {bb12};

    \node[block] (bb6) at (-2,-12.0) {bb6};
    \node[lbl, left=0.15cm of bb6] {back edge};

    \draw[arrow] (bb0) -- (bb1);
    \draw[arrow] (bb1) -- (bb2) node[elbl, midway, left] {iter $<$ 20};
    \draw[arrow] (bb1) -- (bb3) node[elbl, midway, right] {iter $\ge$ 20};
    \draw[arrow] (bb2) -- (bb4) node[elbl, midway, left] {lo $\le$ hi};
    \draw[arrow] (bb2) -- (bb5) node[elbl, midway, right] {lo $>$ hi};
    \draw[arrow] (bb4) -- (bb7) node[elbl, midway, left] {$=$};
    \draw[arrow] (bb4) -- (bb8) node[elbl, midway, right] {$\ne$};
    \draw[arrow] (bb7) -- (bb9);
    \draw[arrow] (bb8) -- (bb10) node[elbl, midway, left] {$<$};
    \draw[arrow] (bb8) -- (bb11) node[elbl, midway, right] {$\ge$};
    \draw[arrow] (bb10) -- (bb12);
    \draw[arrow] (bb11) -- (bb12);
    \draw[arrow] (bb9) -- (bb6);
    \draw[arrow] (bb12) -- (bb6);
    \draw[arrow] (bb5) -- (bb6);
    \draw[arrow] (bb6.east) -- ++(5.5,0) |- (bb1.east);
  \end{tikzpicture}
  \caption{Control-flow graph of \texttt{binary\_search}. Block bb4 contains the critical array access.}
  \label{fig:bs-cfg}
\end{figure}
Key aspects of the TIR representation:

\begin{center}
\small
\begin{tabularx}{0.96\textwidth}{>{\raggedright\arraybackslash}p{0.27\textwidth} >{\raggedright\arraybackslash}X}
\toprule
\textbf{TIR feature} & \textbf{Effect} \\
\midrule
Loop counter & \texttt{iter} is assigned to virtual register \texttt{v8} with an existential type annotation, capturing $0 \le \texttt{v8} < 20$ without fixing a concrete value. \\
Header phi nodes & Mutable variables (\texttt{lo}, \texttt{hi}, \texttt{result}) become explicit loop-carried registers. Useful refinements, such as \texttt{lo} remaining non-negative, may be preserved existentially. \\
Invariant assertion & The loop invariant appears as an \texttt{AssertConstraint} at the entry of bb2. During code generation this becomes a DTAL \texttt{.assert}, allowing the verifier to re-check it independently. \\
\bottomrule
\end{tabularx}
\end{center}

\subsection{DTAL code generation and verification}
The DTAL output preserves the block structure of the TIR, augmented with entry-state declarations, constraint assumptions, and instruction type annotations.

In the current implementation, \texttt{binary\_search} verifies successfully end-to-end. The first snippet below shows the virtual-register DTAL emitted immediately after code generation. The second shows the corresponding physical DTAL after register allocation and x86-specific expansion, which is the form actually consumed by the final verifier.
\subsubsection{Virtual DTAL: midpoint computation and array access}
\begin{verbatim}
.binary_search_bb4:
    .assume v8 < v7
    .assume (v11 >= 0 && v10 < 10)
    .assume v11 <= v10
    sub v16, v10, v11    : {_synth_0: int | _synth_0 == (hi - lo) }
    mov v17, 2           : int(2)
    div v18, v16, v17    : {_synth_1: int | _synth_1 == ((hi - lo) / 2) }
    add v19, v11, v18    : {_synth_2: int |
                            _synth_2 == (lo + ((hi - lo) / 2)) }
    .assert true
    load v20, [v0 + v19] : int
\end{verbatim}
This virtual DTAL is still close to the source-level reasoning. The branch constraints and invariant have already been translated into register-based assumptions, and the midpoint arithmetic is represented symbolically through the synthesised refinement types on \texttt{v16}, \texttt{v18}, and \texttt{v19}. At this level it is easy to see how the array-bounds proof is carried forward: the load of \texttt{v20} depends on the facts $v11 \ge 0$, $v10 < 10$, and $v11 \le v10$, together with the symbolic definition of the midpoint.
\subsubsection{Physical DTAL: verifier-facing code}
\begin{verbatim}
.binary_search_bb2:
    .assert (r0 >= 0 && r4 < 10)
    cmp r0, r4
    setle lr
    mov r3, lr    : bool
    .type r3: bool
    ble .binary_search_bb4
    jmp .binary_search_bb5
.binary_search_bb4:
    mov lr, r4    : int
    mov r7, r0    : int
    sub lr, lr, r7    : {_synth_0: int | _synth_0 == (hi - lo) }
    mov r10, lr    : {_synth_0: int | _synth_0 == (hi - lo) }
    mov r3, 2    : int(2)
    mov r7, r3    : int
    mov lr, r10    : int
    cqo
    idiv r7
    mov r3, lr    : {_synth_1: int | _synth_1 == ((hi - lo) / 2) }
    mov lr, r0    : int
    mov r7, r3    : int
    add lr, lr, r7    : {_synth_2: int | _synth_2 == (lo + ((hi - lo) / 2)) }
    mov r10, lr    : {_synth_2: int | _synth_2 == (lo + ((hi - lo) / 2)) }
    .assert true
    mov lr, r6    : int
    mov r7, r10   : int
    load lr, [lr + r7]    : int
\end{verbatim}
This is the post-register-allocation DTAL that the final verifier checks. The logical structure is the same as in the virtual DTAL, but the values have been assigned to physical registers and division has been expanded into the x86-64-specific \texttt{cqo}/\texttt{idiv} sequence. The verifier re-derives the types of these arithmetic instructions from their operands, proves that the computed midpoint still satisfies the bounds obligation, and checks state coercion on the loop back edge. This illustrates the key trust-model point: register allocation and target-specific expansion remain untrusted, and the verifier validates their output rather than assuming they are correct.
\subsubsection{Back edge and state coercion}
\begin{verbatim}
.binary_search_bb6:
    mov r3, 1    : int
    mov lr, r9   : int
    mov r7, r3   : int
    add lr, lr, r7    : int
    mov r3, lr   : int
    jmp .binary_search_bb1
\end{verbatim}
At the physical back edge, the loop counter has been allocated to \texttt{r9}. The verifier must show that the exit state of bb6 is compatible with the declared entry state of bb1, including the existential bound on the loop counter and the re-established invariant over the registers holding \texttt{lo} and \texttt{hi}. This is the loop's inductive step in its final, post-register-allocation form.

The worked example is deliberately more detailed than the earlier component descriptions because it shows where the project becomes more than a collection of separate checkers. The frontend proves a source-level array access safe, lowering preserves the relevant facts through SSA, code generation re-expresses them as DTAL assumptions and assertions, register allocation rewrites the program into physical registers, and the verifier still reconstructs enough of the argument to accept the final low-level code. That is the core implementation claim of Veritas in miniature: safety evidence is not only generated at the source level, but survives the path to the executable form at which the untrusted compiler is finally checked.
